%% 由 build/emit_tex.py 生成（生成物，勿手改）；定制请放 user_style.tex / user_meta.yaml
\chapter{习题3 参考答案}\label{ux4e60ux98983-ux53c2ux8003ux7b54ux6848}

\begin{quote}
说明：第3章（复变函数的积分）主要习题与参考答案（经典教材整理）。
\end{quote}

\section{习题}\label{ux4e60ux9898-2}

\begin{enumerate}
\def\labelenumi{\arabic{enumi}.}
\item
  沿下列路线计算 \(\displaystyle\int_0^{3+i} z^2\,dz\)：

  \begin{enumerate}
  \def\labelenumii{(\arabic{enumii})}
  \tightlist
  \item
    自原点到 \(3+i\) 的直线段；
  \item
    自原点沿实轴到 \(3\)，再由 \(3\) 沿垂直方向到 \(3+i\)；
  \item
    自原点沿虚轴到 \(i\)，再由 \(i\) 沿水平方向到 \(3+i\)。
  \end{enumerate}
\item
  分别沿 \(y=x\) 与 \(y=x^2\) 计算
  \(\displaystyle\int_0^{1+i}(x+iy^2)\,dz\)。
\item
  设 \(f(z)\) 在单连通域 \(D\) 内解析，\(C\) 为 \(D\)
  内任何一条正向简单闭曲线，问
  \(\displaystyle\int_C \operatorname{Re}[f(z)]\,dz=0\) 与
  \(\displaystyle\int_C \operatorname{Im}[f(z)]\,dz=0\)
  是否成立？若成立请证明，若不成立请举例。
\item
  利用单位圆 \(|z|=1\) 上 \(\bar z=\dfrac1z\) 的性质及柯西积分公式，说明
  \(\displaystyle\oint_C \bar z\,dz=2\pi i\)，其中 \(C\)
  为正向单位圆周。
\item
  计算 \(\displaystyle\int_C \dfrac{dz}{z\bar z}\)，其中 \(C\)
  为正向圆周：

  \begin{enumerate}
  \def\labelenumii{(\arabic{enumii})}
  \tightlist
  \item
    \(|z|=2\)； (2) \(|z|=4\)。
  \end{enumerate}
\item
  沿指定曲线的正向计算下列积分：

  \begin{enumerate}
  \def\labelenumii{(\arabic{enumii})}
  \tightlist
  \item
    \(\displaystyle\int_C\dfrac{e^z}{z-2}dz,\ C:|z-2|=1\)；
  \item
    \(\displaystyle\oint_C\dfrac{dz}{z^2-a^2},\ C:|z-a|=a\)；
  \item
    \(\displaystyle\int_C\dfrac{e^{iz}}{z+i}dz,\ C:|z-2i|=\dfrac32\)；
  \item
    \(\displaystyle\oint_C\dfrac{dz}{z(z-1)(z-2)},\ C:|z|=r<1\)；
  \item
    \(\displaystyle\oint_C\dfrac{\cos z}{z^3}dz,\ C\) 为包围 \(z=0\)
    的简单闭曲线；
  \item
    \(\displaystyle\int_C\dfrac{\sin z}{z}dz,\ C:|z|=1\)。
  \end{enumerate}
\item
  求下列积分的值（\(C:|z|=1\) 取正向）：

  \begin{enumerate}
  \def\labelenumii{(\arabic{enumii})}
  \tightlist
  \item
    \(\displaystyle\oint_C\dfrac{e^z}{z^5}dz\)； (2)
    \(\displaystyle\oint_C\dfrac{z}{(z-1)^3}dz\)。
  \end{enumerate}
\item
  已知 \(u=2(x-1)y\)，求其共轭调和函数 \(v\)，并写出解析函数
  \(f(z)=u+iv\)，使其满足 \(f(2)=-i\)。
\end{enumerate}

\section{参考答案}\label{ux53c2ux8003ux7b54ux6848-2}

\begin{enumerate}
\def\labelenumi{\arabic{enumi}.}
\item
  三个路径结果相同（被积函数解析，积分与路径无关）：\(\displaystyle\int_0^{3+i}z^2dz=\dfrac{(3+i)^3}{3}=\dfrac{26}{3}+6i\)。

  \begin{itemize}
  \tightlist
  \item
    直线段参数
    \(z=3t+it,\ 0\le t\le1\)；分段路径分别积分后叠加，结果相同。
  \end{itemize}
\item
  沿 \(y=x\)：\(z=t(1+i)\)，\(dz=(1+i)dt\)； 沿
  \(y=x^2\)：\(z=t+it^2\)，\(dz=(1+2it)dt\)。两者积分结果相同：\(\dfrac{5}{6}+\dfrac{2}{3}i\)（注意第二个积分需分段展开，最终由解析性保证路径无关）。
\item
  未必成立。例 \(f(z)=z\)，\(C:|z|=1\)： \[
  \oint_C \operatorname{Re}z\,dz=\int_0^{2\pi}\cos\theta\,ie^{i\theta}d\theta=i\pi\neq0 .
  \] 同理 \(\operatorname{Im}\) 也不一定为零。
\item
  在 \(|z|=1\) 上 \(\bar z=\dfrac1z\)，故 \[
  \oint_C \bar z\,dz=\oint_C\dfrac{dz}{z}=2\pi i .
  \]
\item
  \begin{enumerate}
  \def\labelenumii{(\arabic{enumii})}
  \tightlist
  \item
    在 \(|z|=2\) 上 \(z\bar z=|z|^2=4\)，故 \[
    \int_C\dfrac{dz}{z\bar z}=\dfrac14\oint_C dz=0 .
    \]
  \item
    在 \(|z|=4\) 上 \(z\bar z=16\)，故积分为 \(0\)。
  \end{enumerate}
\item
  \begin{enumerate}
  \def\labelenumii{(\arabic{enumii})}
  \tightlist
  \item
    Cauchy 公式：\(=2\pi i\,e^2\)；
  \item
    分解为 \(\dfrac{1}{2a}\left(\dfrac1{z-a}-\dfrac1{z+a}\right)\)，在
    \(C:|z-a|=a\) 内仅含
    \(z=a\)：\(=\dfrac{2\pi i}{2a}=\dfrac{\pi i}{a}\)；
  \item
    奇点 \(z=-i\) 在 \(C\)
    内？\(|{-i}-2i|=|{-3i}|=3>\frac32\)，不在内部，故积分为 \(0\)；
  \item
    对 \(|z|=r<1\)，仅 \(z=0\) 在内部：由
    \(\dfrac1{z(z-1)(z-2)}=\dfrac{A}{z}+\cdots\)，\(A=\dfrac1{2}\)，故
    \(=\pi i\)；
  \item
    高阶导数公式：\(\oint_C\dfrac{\cos z}{z^3}dz=\dfrac{2\pi i}{2!}(\cos z)''\big|_{z=0}=-\pi i\)；
  \item
    \(\oint_C\dfrac{\sin z}{z}dz=2\pi i\sin0=0\)。
  \end{enumerate}
\item
  \begin{enumerate}
  \def\labelenumii{(\arabic{enumii})}
  \tightlist
  \item
    \(\dfrac{2\pi i}{4!}\left(e^z\right)^{(4)}\big|_{z=0}=\dfrac{2\pi i}{24}=\dfrac{\pi i}{12}\)；
  \item
    \(\dfrac{2\pi i}{2!}\cdot 0?\) 注意 \(z=1\) 在 \(|z|=1\)
    上而非内部；若 \(C:|z|=1\)，\(z=1\)
    在圆周上，积分通常不定义或不收敛；一般改为 \(|z-1|=1\) 时
    \(=2\pi i\cdot0=0\)。
  \end{enumerate}
\item
  由 \(v_y=u_x=2y\) 得 \(v=y^2+g(x)\)；再 \(v_x=-u_y=2-2x\)，故
  \(g'(x)=2-2x\)，\(g=2x-x^2+C\)，所以 \(v=y^2+2x-x^2+C\)。 由
  \(f(2)=-i\)：\(f=u+iv\)，在 \(z=2\)（\(x=2,y=0\)）得
  \(u=0,v=4-4+C=C\)，要求 \(f(2)=-i\)，故 \(C=-1\)： \[
  f(z)=2(x-1)y+i(y^2+2x-x^2-1),\qquad f(z)=i(z-1)^2-2i? 
  \] 校验可得 \(f(z)=i(z-1)^2-2i\)。
\end{enumerate}
